% Standalone PCE introduction; the remaining Model chapter stays in model.tex.
\begin{frame}[more-space]{Model\enspace|\enspace Polynomial Chaos Expansion}
  \begin{beamercolorbox}[sep=2mm,wd=\textwidth]{data panel}
    \footnotesize
    A \textbf{Polynomial Chaos Expansion (PCE)} approximates a response as a
    weighted sum of orthonormal polynomial terms.
    \[
      \widehat{T}_{\max}^{\mathrm{ss}}(\mathbf{x})
      = \sum_{|\boldsymbol{\alpha}|\leq p}
        c_{\boldsymbol{\alpha}}\,\Psi_{\boldsymbol{\alpha}}(\boldsymbol{\xi})
    \]
    $c_{\boldsymbol{\alpha}}$: fitted weights; $\Psi_{\boldsymbol{\alpha}}$: polynomial terms;
    $p$: maximum total polynomial degree; $\boldsymbol{\xi}$: scaled inputs.
  \end{beamercolorbox}
  \smallskip
  \begin{columns}[T,onlytextwidth]
    \column{0.48\textwidth}
      \footnotesize
      \textbf{Inputs}\par
      Total CPU power [W]\par
      Inlet air temperature [${}^\circ$C]\par
      Raw fan signal [unit unverified]

    \column{0.49\textwidth}
      \footnotesize
      \textbf{Outputs}\par
      Maximum of the \textbf{18 CPU temperature averages} over the same
      selected steady interval [${}^\circ$C].
  \end{columns}
  \smallskip
  \begin{columns}[c,onlytextwidth]
    \column{0.78\textwidth}
      \begin{beamercolorbox}[sep=2mm,wd=\linewidth]{data panel}
        \footnotesize
        We use \textbf{UQLab} library of ETH Z\"urich for the input models and as our chosen
        PCE framework.\par
        \smallskip
        It provides polynomial bases, regression, UQ and Sobol sensitivity tools
      \end{beamercolorbox}
    \column{0.18\textwidth}
      \centering
      \href{https://www.uqlab.com/}{\includegraphics[width=16mm,height=16mm,keepaspectratio]{images/uqlab.png}}
  \end{columns}
  \datasource{Sources: \href{https://www.uqlab.com/}{UQLab, ETH Z\"urich}}
\end{frame}
